% =====================================================================
% MICCAI / Springer LNCS -- versión en español de partial.tex
% Requiere: llncs.cls + babel (español). Figuras en ./figures/.
% Marcadores a completar antes del envío: [[ ... ]]
% =====================================================================
\documentclass[runningheads]{llncs}

\usepackage[utf8]{inputenc}
\usepackage[T1]{fontenc}
\usepackage[spanish,es-noshorthands]{babel}
\usepackage{graphicx}
\usepackage{booktabs}
\usepackage{amsmath,amssymb}
\usepackage{multirow}
\usepackage{xcolor}
\usepackage{tikz}
\usetikzlibrary{arrows.meta,positioning,fit,backgrounds}
\usepackage[hidelinks]{hyperref}

% Encabezados en español para el entorno abstract / palabras clave de llncs
\renewcommand{\abstractname}{Resumen}
\renewcommand{\keywordname}{\bfseries Palabras clave:}

\newcommand{\todo}[1]{\textcolor{red}{[[#1]]}}

\title{Redes Convolucionales Cuánticas frente a Clásicas para la Segmentación\\
de Subregiones del Tronco Encefálico en la Ataxia Espinocerebelosa Tipo~2:\\
Un Estudio Comparativo con Parámetros Equiparados --- Resultados Parciales}
\titlerunning{Redes Cuánticas vs Clásicas: Segmentación del Tronco Encefálico (SCA2) -- Resultados Parciales}

\author{}
\authorrunning{Resultados Parciales}
\institute{}

\begin{document}
\maketitle

\begin{abstract}
La segmentación precisa de las subregiones del tronco encefálico (bulbo
raquídeo, protuberancia, mesencéfalo) sustenta los biomarcadores de imagen de
la neurodegeneración, en particular la atrofia progresiva de la protuberancia y
el bulbo característica de la ataxia espinocerebelosa tipo~2 (SCA2), un
trastorno cuya mayor prevalencia mundial se registra en Holguín, Cuba.
Presentamos una comparación controlada de nueve arquitecturas de segmentación
sobre una cohorte relevante para la SCA2 adquirida en el Centro de
Neurociencias de Cuba, enfrentando seis U-Nets clásicas (convolucionales y de
tipo transformer) a tres redes híbridas cuántico--clásicas, cada una en los
regímenes 2D (cortes axiales) y 3D (volumétrico) con un presupuesto equiparado
de ${\sim}17$M parámetros. Todos los modelos alcanzan una precisión elevada
(Dice medio $0{,}92$--$0{,}94$). Destacan dos hallazgos. Primero, una ablación
de capacidad ---ampliar el codificador clásico de las redes cuánticas de
$6{,}1$M a $16{,}9$M parámetros sin mejora alguna--- sitúa su techo de
rendimiento en el cuello de botella cuántico, de modo que alcanzan la precisión
de los modelos clásicos con ${\sim}\tfrac{1}{3}$ de los parámetros, a un costo
computacional (simulado) un orden de magnitud mayor. Segundo, la metodología de
evaluación domina las conclusiones en 2D: bajo la métrica convencional
(solo cortes con estructura) todos los modelos 2D son estadísticamente
indistinguibles (Friedman $p{=}0{,}637$), pero bajo una evaluación honesta sobre
el volumen completo se separan con claridad (Friedman $p{<}0{,}0001$), siendo
los modelos híbridos cuánticos los \emph{más robustos} frente a falsos
positivos en cortes vacíos (agrupado $p{=}0{,}031$) y degradándose más algunos
modelos clásicos elaborados. En 3D, ya evaluado sobre volúmenes completos, los
modelos clásicos conservan una pequeña ventaja ($p{=}0{,}008$). Todos los
procedimientos siguieron la Declaración de Helsinki.
\keywords{Segmentación del tronco encefálico \and Ataxia espinocerebelosa \and
Aprendizaje automático cuántico \and U-Net \and Metodología de evaluación.}
\end{abstract}

% =====================================================================
\section{Introducción}
La ataxia espinocerebelosa tipo~2 (SCA2) es un trastorno neurodegenerativo
autosómico dominante causado por una expansión de repeticiones CAG en
\emph{ATXN2}~\cite{pulst}. Su neuropatología se centra en el sistema
olivopontocerebeloso, y la atrofia cuantitativa de la protuberancia y el bulbo
raquídeo constituye un biomarcador de imagen establecido y sensible a la
progresión~\cite{velazquez}. La SCA2 reviste especial importancia en Cuba, donde
la provincia de Holguín presenta la mayor prevalencia del mundo y donde el
Centro de Neurociencias de Cuba mantiene un programa de investigación de larga
trayectoria~\cite{velazquez}. La delineación automática y fiable de las
subregiones del tronco encefálico es un requisito previo para derivar tales
biomarcadores a gran escala; el trazado manual es lento y dependiente del
observador.

Las U-Nets convolucionales profundas~\cite{unet,unet3d} son la herramienta
estándar, pero su huella en parámetros motiva el interés por sesgos inductivos
alternativos. Las redes híbridas cuántico--clásicas se han propuesto como
aproximadores de funciones eficientes en parámetros~\cite{qcnn}, aunque su
utilidad para la segmentación médica volumétrica apenas se ha probado y rara vez
se compara \emph{a igualdad de capacidad}. Planteamos una pregunta
deliberadamente controlada: \textbf{cuando las redes de segmentación clásicas y
cuánticas reciben el mismo presupuesto de parámetros sobre los mismos datos
clínicos, ¿difieren --- y depende la respuesta de cómo las evaluamos?}

\noindent\textbf{Contribuciones.}
(i)~Una comparación reproducible de nueve arquitecturas (seis clásicas, tres
híbridas cuánticas) en 2D y 3D sobre una cohorte de tronco encefálico relevante
para la SCA2.
(ii)~Una ablación de capacidad que sitúa el techo de los modelos cuánticos en el
cuello de botella cuántico, cuantificando una ventaja real de eficiencia en
parámetros (no de precisión) y su costo computacional.
(iii)~Evidencia de que las conclusiones en 2D dependen del protocolo de
evaluación: la evaluación honesta sobre el volumen completo revierte un aparente
resultado de ``ausencia de diferencias'' y revela a los modelos híbridos
cuánticos como los más robustos.

% =====================================================================
\section{Materiales y Métodos}

\subsection{Cohorte, ética y preprocesamiento}
Se adquirieron imágenes de RM cerebral potenciadas en T1 de \todo{$N$} sujetos
(\todo{edad, sexo, composición clínica}) en el Centro de Neurociencias de Cuba.
\textbf{Todos los procedimientos se ajustaron a la Declaración de
Helsinki}~\cite{helsinki}; el estudio fue aprobado por el
\todo{comité de ética, n.\ de aprobación} y todos los participantes otorgaron
su consentimiento informado por escrito. Los volúmenes se procesaron con
corrección de campo de sesgo N4~\cite{n4}, registro afín al espacio MNI y
recorte a un campo de visión de $80\times80\times96$ centrado en el tronco
encefálico; las intensidades se normalizaron por volumen mediante puntuación
$z$ (sin estadísticas entre sujetos, por lo que no hay fuga
entrenamiento/prueba). Las etiquetas expertas definen cuatro clases: fondo,
bulbo raquídeo, protuberancia y mesencéfalo. Los sujetos se dividieron $27/6/6$
(entrenamiento/validación/prueba) a nivel de \emph{sujeto} con una semilla fija;
los seis sujetos de prueba son idénticos para todos los modelos.

\subsection{Arquitecturas}
Evaluamos seis codificadores--decodificadores clásicos ---una \textbf{U-Net} de
doble convolución~\cite{unet,unet3d}, \textbf{CerebNet}~\cite{cerebnet} (bloques
densos competitivos), \textbf{MipaimUNet} (inception multi-rama~\cite{inception}
con conexiones de salto refinadas por atención),
\textbf{SwinUNETR}~\cite{swinunetr,swin} (autoatención por ventanas),
\textbf{ACAPULCO}~\cite{acapulco} (residual con preactivación) y \textbf{MARIN}
(bloques multiescala de preactivación con dilatación dual y atención
CBAM~\cite{cbam})--- y tres modelos híbridos cuántico--clásicos.

\paragraph{Modelos híbridos cuánticos.} \textbf{PennyLane-QCNN}~\cite{pennylane}
y \textbf{Qiskit-QCNN}~\cite{qiskit} son U-Nets clásicas cuyo cuello de botella
se modula mediante un circuito cuántico variacional~\cite{qcnn}: el mapa de
características más profundo $x$ se reduce por promediado global (GAP) a un vector
de canales (se descarta la estructura espacial), se codifica en ángulos sobre
$n_q{=}8$ ($4$) qubits, se transforma con capas de entrelazamiento, se lee como
valores esperados de Pauli-$Z$, se reproyecta y se \emph{suma} a $x$
($x \leftarrow x + q_{\text{feat}}$). La vía cuántica es, por tanto, una
\emph{rama lateral} global y espacialmente uniforme de recalibración de canales
(Fig.~\ref{fig:arch}) --- no un cuello de botella por el que fluya la
representación; las características convolucionales pasan inalteradas. El
\textbf{Hybrid-QCNN} es idéntico a PennyLane-QCNN salvo en la profundidad del
ansatz. En consecuencia, los tres modelos son híbridos; ninguno es puramente
cuántico. Los circuitos se simulan de forma clásica, no se ejecutan en hardware
cuántico.

\begin{figure}[t]
\centering
\begin{tikzpicture}[
  font=\scriptsize,
  box/.style={draw,rounded corners,minimum height=7mm,minimum width=9mm,align=center},
  qbox/.style={box,fill=red!10,draw=red!60},
  cbox/.style={box,fill=blue!6},
  >/.tip={Stealth},every edge/.style={draw,-Stealth}
]
\node[cbox] (in) {Entrada\\$1{\times}80{\times}80$};
\node[cbox,right=6mm of in] (enc) {Codif.\\conv.\\(3 etapas)};
\node[cbox,right=6mm of enc] (bn) {Cuello de\\botella $x$};
\node[circle,draw,inner sep=1pt,right=8mm of bn] (add) {$+$};
\node[cbox,right=8mm of add] (dec) {Decod.\\conv.};
\node[cbox,right=6mm of dec] (head) {conv\\$1{\times}1$};
\node[qbox,above=9mm of add] (qc) {Circuito\\cuántico};
\node[qbox,left=4mm of qc] (gap) {GAP\\$\to n_q$};
\node[qbox,right=4mm of qc] (po) {Lineal\\$\to C$};
\draw (in) edge (enc);
\draw (enc) edge (bn);
\draw (bn) edge (add);
\draw (add) edge (dec);
\draw (dec) edge (head);
\draw[-Stealth] (bn.north) |- (gap.west);
\draw (gap) edge (qc);
\draw (qc) edge (po);
\draw[-Stealth] (po.east) -| (add.north);
\draw[dashed,-Stealth] (enc.south) to[out=-35,in=-145] (dec.south);
\node[font=\tiny,below=4mm of add,red!70] {conexiones de salto (discontinuas)};
\begin{scope}[on background layer]
\node[draw=red!40,rounded corners,fill=red!4,fit=(gap)(qc)(po),inner sep=2mm,
  label={[red!70,font=\tiny]above:rama lateral cuántica aditiva}] {};
\end{scope}
\end{tikzpicture}
\caption{Arquitectura híbrida cuántico--clásica. El circuito cuántico opera
sobre un vector de canales obtenido por promediado global (GAP) y su salida se
difunde y se \emph{suma} al cuello de botella convolucional $x$; es una rama
lateral de recalibración de canales, no un verdadero cuello de botella.}
\label{fig:arch}
\end{figure}

\paragraph{Equiparación de parámetros.} Los modelos clásicos autoseleccionan
anchos de canal específicos por dimensión para mantener ${\sim}17$M parámetros
en 2D y 3D. Los modelos cuánticos/híbridos usan un codificador fijo de $3$
etapas $[80,160,320]$ que da ${\sim}6{,}1$M en 2D pero ${\sim}17{,}6$M en 3D
(los núcleos 3D cuestan ${\sim}3\times$). Los modelos cuánticos 2D nativos, por
tanto, \emph{no} están equiparados; añadimos una segunda ejecución 2D con un
codificador de $4$ etapas $[80,160,320,480]\,({=}16{,}9$M, las variantes
``17M'').

\subsection{Entrenamiento, métricas y hardware}
Los modelos se entrenaron con Adam~\cite{adam} (tasa~$10^{-3}$, programación
coseno, calentamiento de $5$ épocas), una pérdida combinada de entropía cruzada
y Dice suave~\cite{vnet} con ponderación de clases por frecuencia inversa,
precisión mixta, recorte de gradiente y parada temprana (paciencia~$15$,
mínimo $250$ épocas) sobre una única NVIDIA RTX~3090. Reportamos el Dice por
clase de primer plano y el Dice medio de primer plano. La unidad estadística
primaria es el Dice medio por sujeto sobre los seis sujetos de prueba; dado el
$n$ reducido reportamos intervalos de confianza del $95\%$ por remuestreo
(bootstrap) de las diferencias pareadas y la $d$ de Cohen pareada como evidencia
primaria, con la prueba ómnibus de Friedman y pruebas de Wilcoxon corregidas por
Holm~\cite{nonparam} como exploratorias.

\paragraph{Protocolo de evaluación.} Para el entrenamiento 2D se descartan los
cortes axiales con $<20$ vóxeles de primer plano. Aplicar este filtro al
conjunto de prueba puntúa \emph{solo} los cortes con tronco encefálico, sin
penalizar nunca los falsos positivos en cortes vacíos, dando un Dice optimista.
Por ello reportamos además una \textbf{evaluación honesta sobre el volumen
completo} en 2D sobre los $96$ cortes axiales de cada sujeto. Los modelos 3D ya
evalúan volúmenes completos; en consecuencia, las cifras 2D y 3D no son
directamente comparables. Los experimentos se ejecutaron en una estación con
$2{\times}$RTX~3090 (24\,GB) (una por ejecución), Python~3.12, PyTorch~2.6,
PennyLane~0.44 y Qiskit~2.4.

% =====================================================================
\section{Resultados}
Todos los valores de Dice son la media por sujeto~$\pm$~d.e.\ sobre los seis
sujetos de prueba.

\subsection{Precisión a igualdad de capacidad}
Todas las arquitecturas alcanzaron una precisión elevada
(Tabla~\ref{tab:main}). Sobre cortes 2D con estructura, los nueve modelos a
igualdad de capacidad fueron estadísticamente indistinguibles
(Friedman $\chi^2{=}6{,}09$, $p{=}0{,}637$); el mejor fue un modelo cuántico
(Qiskit-QCNN). En 3D los modelos difirieron ($\chi^2{=}20{,}58$,
$p{=}0{,}008$) y el mejor fue clásico (MipaimUNet).

\begin{table}[t]
\centering
\caption{Dice medio por sujeto $\pm$ d.e.\ ($n{=}6$). Q=cuántico, H=híbrido,
C=clásico. El 2D se muestra sobre cortes con estructura al presupuesto
equiparado de $16{,}9$M.}
\label{tab:main}
\small
\setlength{\tabcolsep}{4pt}
\begin{tabular}{lcc@{\hskip 1em}lcc}
\toprule
\multicolumn{3}{c}{\textbf{2D (cortes axiales)}} &
\multicolumn{3}{c}{\textbf{3D (volúmenes)}}\\
\cmidrule(r){1-3}\cmidrule(l){4-6}
Modelo & Par. & Dice & Modelo & Par. & Dice\\
\midrule
Qiskit-QCNN (Q)    & 16,9M & $.932{\pm}.007$ & MipaimUNet (C)     & 17,8M & $.940{\pm}.013$\\
MARIN (C)          & 17,4M & $.932{\pm}.015$ & ACAPULCO (C)       & 17,1M & $.937{\pm}.017$\\
Hybrid-QCNN (H)    & 16,9M & $.932{\pm}.011$ & PennyLane-QCNN (Q) & 17,6M & $.935{\pm}.016$\\
MipaimUNet (C)     & 17,3M & $.931{\pm}.015$ & MARIN (C)          & 17,1M & $.934{\pm}.019$\\
PennyLane-QCNN (Q) & 16,9M & $.930{\pm}.013$ & SwinUNETR (C)      & 17,1M & $.932{\pm}.020$\\
SwinUNETR (C)      & 16,8M & $.929{\pm}.014$ & Hybrid-QCNN (H)    & 17,6M & $.931{\pm}.022$\\
ClassicalCNN (C)   & 17,0M & $.929{\pm}.013$ & CerebNet (C)       & 18,0M & $.930{\pm}.018$\\
ACAPULCO (C)       & 17,4M & $.928{\pm}.017$ & Qiskit-QCNN (Q)    & 17,6M & $.925{\pm}.032$\\
CerebNet (C)       & 16,7M & $.927{\pm}.016$ & ClassicalCNN (C)   & 17,0M & $.920{\pm}.037$\\
\bottomrule
\end{tabular}
\end{table}

\subsection{El cuello de botella cuántico es el techo}
Ampliar el codificador clásico de los modelos cuánticos/híbridos 2D de
$6{,}1$M a $16{,}9$M parámetros no produjo \emph{ninguna} mejora
(PennyLane $\Delta{=}{-}0{,}0023$, IC~95\% $[-0{,}0046,-0{,}0005]$;
Qiskit $\Delta{=}{-}0{,}0007$; Hybrid $\Delta{=}{+}0{,}0019$; ninguno sobrevive
a la corrección de Holm). El rendimiento está limitado por el cuello de botella
de $8$ qubits, no por la red troncal --- de modo que los modelos cuánticos
alcanzan el Dice de los clásicos con ${\sim}\tfrac13$ de los parámetros. El
costo es computacional: las ejecuciones cuánticas 2D tardaron $8{,}9$--$13{,}6$\,h
frente a $0{,}3$--$1{,}0$\,h de los modelos clásicos; en 3D la brecha se reduce
(cuántico $28$--$35$\,min) al dominar el costo la red clásica de mayor tamaño
(Fig.~\ref{fig:cost}).

\subsection{La evaluación honesta sobre el volumen completo revierte el
resultado en 2D}
Repuntuar los mismos modelos 2D sobre los $96$ cortes de cada sujeto cambia el
panorama sustancialmente (Tabla~\ref{tab:c2}, Fig.~\ref{fig:c2}). La prueba de
Friedman salta de $p{=}0{,}637$ a $\chi^2{=}39{,}5$, $p{<}0{,}0001$, y la
dispersión del Dice medio se amplía de ${\approx}0{,}006$ a $0{,}059$. El filtro
había enmascarado grandes diferencias en el comportamiento de falsos positivos
en cortes vacíos: MARIN ---co-líder en cortes filtrados--- cae al \emph{último}
puesto ($-0{,}087$), mientras que los modelos híbridos cuánticos son los más
robustos. Agrupado, cuántico/híbrido vs clásico: $0{,}894$ vs $0{,}878$,
$\Delta{=}{+}0{,}016$, Wilcoxon $p{=}0{,}031$. Los modelos cuánticos estrechos de
$6{,}1$M muestran las menores caídas de todos.

\begin{table}[t]
\centering
\caption{Dice 2D: evaluación convencional (solo cortes con estructura) frente a
volumen completo honesto (media por sujeto, $n{=}6$). $\Delta$ = volumen
completo $-$ filtrado.}
\label{tab:c2}
\setlength{\tabcolsep}{6pt}
\begin{tabular}{llccc}
\toprule
Modelo & Tipo & Filtrado & Vol.\ completo & $\Delta$\\
\midrule
Qiskit-QCNN (17M)    & Q & 0,932 & \textbf{0,905} & $-0,028$\\
CerebNet             & C & 0,927 & 0,899 & $-0,028$\\
ClassicalCNN         & C & 0,929 & 0,898 & $-0,031$\\
Hybrid-QCNN (17M)    & H & 0,932 & 0,898 & $-0,034$\\
SwinUNETR            & C & 0,929 & 0,883 & $-0,046$\\
PennyLane-QCNN (17M) & Q & 0,930 & 0,880 & $-0,050$\\
MipaimUNet           & C & 0,931 & 0,875 & $-0,056$\\
ACAPULCO             & C & 0,928 & 0,866 & $-0,062$\\
MARIN                & C & 0,932 & \textbf{0,845} & $-0,087$\\
\midrule
PennyLane-QCNN (6M)  & Q & 0,932 & 0,914 & $-0,018$\\
Qiskit-QCNN (6M)     & Q & 0,933 & 0,912 & $-0,022$\\
Hybrid-QCNN (6M)     & H & 0,930 & 0,899 & $-0,031$\\
\bottomrule
\end{tabular}
\end{table}

\begin{figure}[t]
\centering
\includegraphics[width=0.78\textwidth]{figures/fig_fullvol_vs_fg_2d_es.png}
\caption{Dice 2D por sujeto bajo evaluación filtrada (gris) frente a volumen
completo honesto (color: rojo = cuántico/híbrido, azul = clásico), ordenado por
la puntuación de volumen completo. La métrica filtrada comprime todos los
modelos cerca de $0{,}93$; la evaluación honesta los separa, siendo los
cuánticos/híbridos los más robustos.}
\label{fig:c2}
\end{figure}

\begin{figure}[t]
\centering
\includegraphics[width=0.62\textwidth]{figures/fig_cost_2d_es.png}
\caption{Precisión 2D frente al tiempo de entrenamiento (escala log). Los
modelos cuánticos ocupan la región de alto costo con precisión comparable, lo
que refleja la simulación clásica de los circuitos.}
\label{fig:cost}
\end{figure}

% =====================================================================
\section{Discusión}
A igualdad de capacidad, las U-Nets híbridas cuántico--clásicas igualan a las
redes clásicas en la segmentación de subregiones del tronco encefálico en 2D y
quedan ligeramente por detrás en 3D. La ablación muestra que su ventaja es de
\emph{eficiencia en parámetros, no de precisión}: un cuello de botella de $8$
qubits sustituye a varios millones de parámetros convolucionales sin pérdida de
Dice, pero también limita la precisión alcanzable; hoy esto se compensa con el
costo de la simulación de vector de estado.

El hallazgo más relevante es metodológico. El protocolo 2D convencional
(solo cortes con estructura) hacía que todos los modelos parecieran idénticos,
mientras que la puntuación honesta sobre el volumen completo ---que penaliza los
falsos positivos en cortes sin tronco encefálico--- los separó en un orden de
magnitud y reordenó el ranking. Los modelos híbridos cuánticos, y en especial
las variantes estrechas de $6{,}1$M, resultaron los más robustos, posiblemente
porque la rama cuántica global aditiva actúa como una compuerta conservadora de
canales que suprime activaciones espurias ante entradas vacías. Esto sugiere que
los \textbf{benchmarks de segmentación del tronco encefálico deben reportarse
sobre volúmenes completos}, no sobre cortes filtrados, para ser clínicamente
significativos. Clínicamente, todos los modelos produjeron segmentaciones lo
bastante precisas (Dice de bulbo/protuberancia ${\sim}0{,}94$) para sustentar
las mediciones de atrofia de la protuberancia y el bulbo relevantes para la
SCA2; para su despliegue, una U-Net 3D clásica y rápida sigue siendo la opción
pragmática.

\noindent\textbf{Limitaciones.} El conjunto de prueba tiene seis sujetos, lo que
limita la potencia estadística; las diferencias se reportan como tamaños del
efecto con intervalos por remuestreo y se interpretan como exploratorias. Cada
configuración usó una única semilla; los datos son de un solo centro; los
modelos cuánticos se simularon de forma clásica. Una auditoría del código también
identificó un término de Dice ponderado por clases mal escalado (ya corregido);
los modelos reportados se entrenaron con la pérdida original, por lo que la
comparación \emph{relativa} es internamente consistente pero los valores
absolutos pueden variar tras reentrenar. Cohortes multicéntricas, entrenamiento
con múltiples semillas, reentrenamiento con volumen completo y validación
prospectiva de biomarcadores son los siguientes pasos naturales.

% =====================================================================
\section{Conclusión}
Sobre una cohorte de tronco encefálico relevante para la SCA2, las redes de
segmentación cuánticas y clásicas con parámetros equiparados tienen una
precisión comparable, ofreciendo los modelos cuánticos eficiencia en parámetros
a un costo computacional sustancial. De forma crucial, las conclusiones en 2D
dependen del protocolo de evaluación: la puntuación honesta sobre el volumen
completo revela a los modelos híbridos cuánticos como los más robustos frente a
falsos positivos en cortes vacíos. Los hallazgos respaldan la investigación
continua y cuidadosamente controlada de la segmentación híbrida
cuántico--clásica y de estándares de evaluación sobre volúmenes completos.

\subsubsection*{Agradecimientos.} \todo{Financiación y agradecimientos de
recolección de datos.}

% =====================================================================
\begin{thebibliography}{99}
\bibitem{unet} Ronneberger, O., Fischer, P., Brox, T.: U-Net: Convolutional networks for biomedical image segmentation. In: MICCAI (2015)
\bibitem{unet3d} \c{C}i\c{c}ek, \"O., et al.: 3D U-Net: learning dense volumetric segmentation from sparse annotation. In: MICCAI (2016)
\bibitem{vnet} Milletari, F., Navab, N., Ahmadi, S.A.: V-Net: fully convolutional neural networks for volumetric medical image segmentation. In: 3DV (2016)
\bibitem{cerebnet} Faber, J., et al.: CerebNet: a fast and reliable deep-learning pipeline for detailed cerebellum sub-segmentation. NeuroImage (2022)
\bibitem{acapulco} Han, S., Carass, A., et al.: Automatic cerebellum anatomical parcellation using U-Net with locally constrained optimization (ACAPULCO). NeuroImage (2020)
\bibitem{swinunetr} Hatamizadeh, A., et al.: Swin UNETR: Swin transformers for semantic segmentation of brain tumors in MRI images. In: MICCAI BrainLes Workshop (2022)
\bibitem{swin} Liu, Z., et al.: Swin Transformer: hierarchical vision transformer using shifted windows. In: ICCV (2021)
\bibitem{inception} Szegedy, C., et al.: Going deeper with convolutions. In: CVPR (2015)
\bibitem{cbam} Woo, S., Park, J., Lee, J.Y., Kweon, I.S.: CBAM: convolutional block attention module. In: ECCV (2018)
\bibitem{qcnn} Cong, I., Choi, S., Lukin, M.D.: Quantum convolutional neural networks. Nature Physics 15, 1273--1278 (2019)
\bibitem{pennylane} Bergholm, V., et al.: PennyLane: automatic differentiation of hybrid quantum-classical computations. arXiv:1811.04968 (2018)
\bibitem{qiskit} Qiskit contributors: Qiskit: an open-source framework for quantum computing (2023)
\bibitem{n4} Tustison, N.J., et al.: N4ITK: improved N3 bias correction. IEEE TMI 29(6), 1310--1320 (2010)
\bibitem{adam} Kingma, D.P., Ba, J.: Adam: a method for stochastic optimization. In: ICLR (2015)
\bibitem{pulst} Pulst, S.M., et al.: Moderate expansion of a normally biallelic trinucleotide repeat in spinocerebellar ataxia type 2. Nature Genetics 14, 269--276 (1996)
\bibitem{velazquez} Velazquez-Perez, L., Rodriguez-Labrada, R., Fernandez-Ruiz, J.: Spinocerebellar ataxia type 2: clinicogenetic aspects, mechanistic insights, and management. \todo{verificar revista/año}
\bibitem{helsinki} Asociación Médica Mundial: Declaración de Helsinki: principios éticos para las investigaciones médicas en seres humanos. JAMA 310(20), 2191--2194 (2013)
\bibitem{nonparam} Friedman, M.: The use of ranks to avoid the assumption of normality implicit in the analysis of variance. JASA 32, 675--701 (1937); Wilcoxon, F.: Individual comparisons by ranking methods. Biometrics Bulletin 1, 80--83 (1945)
\end{thebibliography}

\end{document}
